\section{Introduction}
\label{sec:introduction}

Minimum decycling sets in directed de Bruijn graphs are classical
combinatorial objects and underlie several finite-window sequence-sampling
constructions.  Their ordinary minimum is governed by rotation cycles and
unavoidable sets \cite{Mykkeltveit1972,ChamparnaudHanselPerrin2004}.  More
recent work made the move structure of minimum decycling sets explicit
\cite{MarcaisDeBlasioKingsford2024}.  Reverse-complement symmetry imposes a
different constraint: whenever a word is chosen, its reverse complement must
also be chosen.  This is the natural strand symmetry for DNA and was posed
directly in the k-nonical setting \cite{MarcaisElderKingsford2024}.

In Section~5.1 of \cite{MarcaisElderKingsford2024}, the authors identify
the lack of known algorithms for symmetric MDSs; Section~6 also notes the
unknown symmetric minimum size. We answer the following existence part of
that problem: for which word lengths can literal reverse-complement
symmetry attain the ordinary minimum? The answer below classifies every
fixed-point-free complemented alphabet. It does not determine the exact
odd-order constrained minimum or solve the additional path-length
optimization in that work.

\begin{theorem}[Parity theorem]
\label{thm:parity}
Let $\A$ be a finite nonempty alphabet with a fixed-point-free involutive
complement, and put $q=|\A|$.  For every integer $k\geq2$, the directed de
Bruijn graph $B(q,k)$ has a reverse-complement-invariant ordinary minimum
decycling set if and only if $k$ is even.
\end{theorem}

In particular, Theorem~\ref{thm:parity} applies to binary words and to DNA
words under $A\leftrightarrow T$ and $C\leftrightarrow G$.  The case $k=1$
is deliberately outside the theorem: every loop must be deleted, so the
unique ordinary minimum is the whole alphabet and is invariant even though
$k$ is odd.

The two directions use different mechanisms.  At odd order, a primitive
pure-rotation cycle has two reverse-complement-fixed lower endpoints.  One
deletion from that cycle cannot destroy both endpoint-to-endpoint directions,
and reverse complement supplies the opposing residual path. Counting this
obstruction over disjoint alphabet pairs gives at least $q$ additional
deletions. At even order,
we assign distinct nonzero real weights $s(a)$ satisfying
$s(\bar a)=-s(a)$.  A weighted cyclotomic version of the classical selector
then has an exact reverse-complement defect equal to the gradient of an open
half-plane indicator.  Legal full-alphabet source firings remove that defect
monotonically.

The proof yields two complementary constructions.  The first is a finite
descent and exposes the move-theoretic mechanism; every allowed monotone
scheduler uses an explicitly determined number of firings.  The second cuts
the same self-dual residual poset at a spectral midpoint and gives a direct
word-membership rule.  When the weight table is integer-valued, exact
real-algebraic sign determination makes the latter a polynomial-bit
algorithm.  Neither statement supplies a bound on
the longest residual path or a biological performance guarantee.

\paragraph{Relation to prior terminology.}
The phrase ``symmetric decycling set'' has also been used for the ordinary
Mykkeltveit construction based on opposite geometric half-planes
\cite{PellowEtAl2023}.  We use the unambiguous phrase
\emph{reverse-complement-invariant} throughout.

\paragraph{Evidence discipline.}
The all-order proof occupies Sections~\ref{sec:odd}--\ref{sec:parity-proof}.
Section~\ref{sec:validation} records an independent exact implementation of
the binary specialization only as a sign, convention, and mutation check.
No finite census is a premise of Theorem~\ref{thm:parity}.

\paragraph{Organization.}
Section~\ref{sec:normalization} fixes the graph, rotation, and gradient
conventions.  Section~\ref{sec:odd} proves the odd obstruction.  Sections
\ref{sec:seed}--\ref{sec:even} construct and symmetrize the weighted seed.
Section~\ref{sec:direct} gives the direct spectral selector, and
Section~\ref{sec:parity-proof} assembles the parity theorem and its DNA
corollary.  The appendices collect convention maps, technical details, and
the finite binary assurance table.

\begin{boundary}
The mathematical theorem is alphabet-general under a fixed-point-free
complement.  We do not claim an exact odd-order symmetry premium, an optimal
firing path, a quantitative window guarantee, biological performance, or
blanket novelty or priority over unexamined sources.
\end{boundary}
